\chapter{Perancangan Sistem}

\section{Gambaran Umum Sistem}

Sistem yang dirancang pada penelitian ini merupakan sebuah \textit{pipeline} analisis \textit{process mining} yang memanfaatkan \textit{generative AI} untuk melakukan abstraksi semantik label aktivitas pemeliharaan sebelum model proses dibangun. Titik berangkat perancangan adalah persoalan nyata pada data operasional pemeliharaan pembangkit listrik, yaitu deskripsi pekerjaan pada \textit{job card} yang ditulis secara bebas oleh teknisi sehingga satu jenis pekerjaan dapat muncul dalam puluhan variasi penulisan. Variasi tersebut membuat model proses yang dihasilkan dari \textit{raw labels} menjadi sangat rumit dan sukar dibaca.

Untuk mengatasi hal itu, sistem dirancang dalam tiga tahap inti, sebagaimana digambarkan pada Gambar~\ref{fig:arsitektur_sistem}. Tahap pertama, \textit{preprocessing}, membersihkan \textit{raw data}, memisahkan \textit{compound activity} berbasis aturan, serta menormalkan \textit{Case ID} dan \textit{timestamp} hingga terbentuk \textit{event log} dasar yang siap dianotasi. Tahap kedua, \textit{log construction} dan \textit{process discovery}, menyusun dua versi \textit{event log} dari basis yang sama, yaitu \textit{log as-is} berlabel aktivitas asli dan \textit{log enriched} berlabel kategori tindakan hasil anotasi \textit{Large Language Model} (LLM), lalu membangun model proses dari masing-masing log. Tahap ketiga, \textit{evaluasi}, membandingkan kualitas kedua model melalui \textit{conformance checking}. Bab ini hanya memuat rancangan konseptual; realisasi teknis dan keluarannya dibahas pada Bab~IV.

\begin{figure}[H]
    \centering
    \includegraphics[width=\textwidth]{"bab3/FlowChart Perancangan Sistem-Gambaran Sistem Revamp.png"}
    \caption{Diagram Arsitektur Sistem}
    \label{fig:arsitektur_sistem}
\end{figure}

\section{Rancangan Pipeline}

\textit{Pipeline} disusun secara modular agar tiap tahap dapat diuji dan diganti tanpa mengganggu tahap lain. Modul pertama, \textit{Data Ingestion and Preprocessing}, bertanggung jawab membaca \textit{raw data}, menstandarkan kolom, memperbaiki \textit{Case ID} yang hilang, melengkapi \textit{timestamp} yang kosong, serta memisahkan \textit{compound activity} menjadi \textit{event} tunggal berbasis aturan, hingga terbentuk \textit{event log} dasar yang bersih dan siap dianotasi. Modul kedua, \textit{AI Annotation Builder}, mengirim konteks tiap \textit{event} ke model bahasa, lalu menerima kembali kategori tindakan (\texttt{action\_type}) beserta atribut pendukungnya tanpa mengubah jumlah \textit{event}, sehingga terbentuk \textit{event log} yang diperkaya. Modul ketiga, \textit{Process Model Discovery and Conformance Checking}, membangun model proses berupa \textit{Petri net} dari dua versi log, yaitu \textit{log as-is} dan \textit{log enriched}, lalu mengukur kualitas keduanya melalui empat dimensi \textit{conformance}. Modul keempat, \textit{Report Generation and Visualization}, merangkum temuan ke dalam tabel, grafik, dan laporan.

Tabel~\ref{tab:rancangan-modul-pipeline} merangkum fungsi serta masukan dan keluaran tiap modul pada tataran konseptual. Wujud teknisnya, termasuk nama berkas, ukuran \textit{event log}, dan nilai metrik yang sebenarnya, merupakan keluaran eksekusi yang dilaporkan pada Bab~IV.

\begin{table}[H]
\centering
\caption{Rancangan Modul Pipeline}
\label{tab:rancangan-modul-pipeline}
\footnotesize
\begin{tabularx}{\textwidth}{p{3cm} p{3.2cm} X X}
\toprule
\textbf{Modul} & \textbf{Fungsi} & \textbf{Masukan (konseptual)} & \textbf{Keluaran (konseptual)} \\
\midrule
\textit{Data Ingestion \& Preprocessing} &
Membersihkan data, menormalkan \textit{Case ID} dan \textit{timestamp}, serta memisahkan \textit{compound activity} berbasis aturan &
\textit{Raw data} pemeliharaan (WBS dan \textit{job card}) &
\textit{Event log} dasar yang bersih dan siap dianotasi \\
\addlinespace
\textit{AI Annotation Builder} &
Memberi kategori \texttt{action\_type} dan atribut semantik pada tiap \textit{event} &
\textit{Event log} dasar beserta konteks \textit{event} &
\textit{Event log} diperkaya berlabel \texttt{action\_type} \\
\addlinespace
\textit{Process Discovery \& Conformance} &
Menemukan model proses dari dua versi log dan mengukur kualitasnya &
\textit{Event log as-is} dan \textit{event log enriched} &
Model proses \textit{as-is} dan \textit{enriched} beserta nilai \textit{conformance} \\
\addlinespace
\textit{Report Generation \& Visualization} &
Merangkum temuan menjadi laporan dan visualisasi &
Model proses dan hasil \textit{conformance} &
Laporan dan visualisasi temuan \\
\bottomrule
\end{tabularx}
\end{table}

\section{Perancangan Data Masukan}

\subsection{Jenis dan Sumber Data}

Data yang digunakan berasal dari kegiatan \textit{overhaul} pemeliharaan pada PLTD Titikuning. Sumber utamanya adalah dokumen \textit{Work Breakdown Structure} (WBS) dan \textit{job card} yang mencatat rangkaian pekerjaan pada setiap unit dan komponen. Setelah diekspor menjadi berkas \texttt{event\_log-time.csv}, \textit{raw data} memuat catatan aktivitas pemeliharaan yang masing-masing dilengkapi \textit{timestamp}. Jumlah baris aktual hasil ekspor dilaporkan pada Bab~IV.

\subsection{Struktur \textit{Raw Data}}

\textit{Raw data} tersusun atas sejumlah kolom yang menggambarkan \textit{Case ID}, deskripsi aktivitas, dan \textit{timestamp}. Struktur ini menjadi dasar pemetaan menuju \textit{event log} yang sesuai dengan kebutuhan \textit{process mining}.

\begin{table}[H]
\centering
\caption{Struktur Kolom Raw Data}
\label{tab:struktur-data-mentah}
\footnotesize
\begin{tabularx}{\textwidth}{p{3.5cm} p{2cm} X}
\toprule
\textbf{Nama Kolom} & \textbf{Tipe Data} & \textbf{Deskripsi} \\
\midrule
\texttt{Task ID} & String & Penanda unik \textit{case}/pekerjaan (contoh: TASK-001) \\
\texttt{Activity Description} & String & \textit{Raw activity description} (contoh: ``Pembersihan \& Pengecekan Filter'') \\
\texttt{Actual Start} & Datetime & Waktu mulai pekerjaan (contoh: 2024-01-01 08:00) \\
\texttt{Actual Finish} & Datetime & Waktu selesai pekerjaan (contoh: 2024-01-01 10:00) \\
\bottomrule
\end{tabularx}
\end{table}

\subsection{Pemetaan ke \textit{Event Log}}

Agar dapat diolah oleh teknik \textit{process mining}, \textit{raw data} dipetakan ke dalam tiga atribut wajib \textit{event log}. \textit{Case ID} dipetakan ke \texttt{case:concept:name}, nama aktivitas dipetakan ke \texttt{concept:name} (\textit{Activity}), dan \textit{timestamp} dipetakan ke \texttt{time:timestamp}. Pemetaan ini memungkinkan tiap pekerjaan diurutkan secara kronologis dalam satu \textit{case} pemeliharaan sehingga \textit{traces} proses dapat direkonstruksi.

\section{Preprocessing}

Tahap \textit{preprocessing} dirancang untuk mengubah \textit{raw data} menjadi \textit{event log} yang bersih dan konsisten. Terdapat empat langkah utama yang dijalankan secara berurutan.

Langkah pertama adalah resolusi \textit{Case ID}. Sebagian catatan tidak memiliki \textit{Case ID} yang valid, sehingga sistem membangkitkan \textit{Case ID} otomatis untuk mengelompokkannya.

Langkah kedua adalah \textit{noise filtering}. Baris yang tidak relevan atau tidak lengkap dibersihkan agar tidak mengganggu pembentukan model.

Langkah ketiga adalah imputasi \textit{timestamp}. Sebagian \textit{event} tidak memiliki \textit{timestamp} yang lengkap, sehingga nilai waktunya diperkirakan berdasarkan urutan dan konteks pekerjaan.

Langkah keempat adalah \textit{event splitting}, yaitu pemisahan \textit{compound activity}. Banyak deskripsi pekerjaan menggabungkan beberapa tindakan dalam satu baris, misalnya pembersihan dan pemeriksaan yang ditulis bersamaan, sehingga satu catatan sebenarnya memuat lebih dari satu \textit{event}. Berdasarkan kumpulan kata kunci tindakan, sistem mengenali baris semacam itu lalu memisahkannya menjadi \textit{event} tunggal. Setiap \textit{event} hasil pemisahan ditandai melalui kolom \texttt{is\_split} dan \texttt{split\_of}.

Hasil tahap ini berupa dua \textit{event log}, yaitu \textit{event log as-is} (\texttt{log\_raw}) berlabel aktivitas mentah dan \textit{event log enriched} (\texttt{log\_enriched}) setelah \textit{event splitting}. \textit{Event log enriched} menjadi masukan bagi tahap anotasi pada Subbab~\ref{subsec:annotation-builder}. Ukuran \textit{event log}, jumlah \textit{cases}, serta jumlah imputasi dan pemisahan aktual dilaporkan pada Bab~IV.

\section{Annotation Builder Berbasis LLM}
\label{subsec:annotation-builder}

Tahap anotasi bekerja pada \textit{event log enriched} yang telah memuat \textit{events} hasil \textit{event splitting} pada tahap \textit{preprocessing}. Pemisahan tersebut dilakukan secara aturan; tahap ini berfokus pada abstraksi semantik, yaitu penetapan kategori \texttt{action\_type} beserta atribut pendukung pada tiap \textit{event}, tanpa menambah atau mengurangi jumlah \textit{events}.

\subsection{Taxonomy Boundary}

Inti dari abstraksi semantik adalah penyeragaman tindakan ke dalam sekumpulan kategori baku. Penelitian ini menetapkan dua belas kategori \texttt{action\_type} sebagai batas taksonomi yang menjadi acuan model bahasa. Batas taksonomi diberikan sebagai aturan, sedangkan penentuan kategori untuk tiap \textit{event} diserahkan sepenuhnya kepada model bahasa. Kedua belas kategori tersebut tercantum pada Tabel~\ref{tab:taxonomy_boundary}.

\begin{table}[H]
\centering
\caption{Kategori Action Type (Taxonomy Boundary)}
\label{tab:taxonomy_boundary}
\small
\begin{tabularx}{\textwidth}{c l X}
\toprule
\textbf{No.} & \textbf{Kategori} & \textbf{Cakupan Tindakan} \\
\midrule
1 & SAFETY\_INDUCTION & Pengarahan dan persiapan keselamatan kerja \\
2 & CLEANING & Pembersihan komponen atau area kerja \\
3 & INSPECTION & Pemeriksaan dan pengecekan kondisi \\
4 & MEASUREMENT & Pengukuran parameter teknis dan kalibrasi \\
5 & REPLACEMENT & Penggantian komponen \\
6 & REPAIR & Perbaikan komponen yang rusak \\
7 & DISASSEMBLY & Pelepasan atau pembongkaran komponen \\
8 & ASSEMBLY & Pemasangan kembali komponen \\
9 & TESTING & Pengujian fungsi setelah pekerjaan \\
10 & ADJUSTMENT & Penyetelan dan penyesuaian \\
11 & LUBRICATION & Pelumasan dan penggantian pelumas \\
12 & OTHER & Tindakan lain di luar kategori utama \\
\bottomrule
\end{tabularx}
\end{table}

\subsection{Desain Payload dan Prompt}

Untuk setiap \textit{case}, sistem menyusun \textit{payload} yang berisi \textit{raw context} dari \textit{event}, posisinya dalam urutan \textit{case}, aktivitas sebelum dan sesudahnya, serta konteks WBS terkait. \textit{Payload} dikirim ke model \texttt{gemini-3.1-flash-lite} dengan \textit{temperature} 0,1 agar keluaran bersifat deterministik dan konsisten.

Perancangan \textit{prompt} mengikuti prinsip \textit{taxonomy boundary}: daftar dua belas kategori \texttt{action\_type} disertakan sebagai batas label yang valid, bukan sebagai aturan pencocokan kata kunci. Instruksi sistem meminta model menetapkan satu \texttt{action\_type} per \textit{event} berdasarkan makna deskripsi, tanpa membuat label di luar taksonomi. Keluaran diminta dalam format terstruktur yang mencakup \texttt{action\_type}, nama komponen, subsistem, area WBS, skor \textit{confidence}, dan \textit{semantic basis} pemilihan kategori. \textit{Semantic basis} memuat alasan domain atas pemilihan kategori, bukan rujukan ke aturan atau kata kunci internal sistem.

Untuk menjaga kualitas anotasi, \textit{event} dengan \textit{confidence} di bawah 0,45 diklasifikasikan anomali dan ditinjau ulang. Anotasi berbasis aturan hanya digunakan apabila layanan model bahasa tidak tersedia. Uji \textit{prompt} dilakukan pada sampel data sebelum anotasi penuh dijalankan, guna memastikan pemetaan kategori masuk akal pada konteks pemeliharaan PLTD.

\section{Skenario Process Discovery}

Tahap \textit{process discovery} dilakukan melalui dua skenario yang dibandingkan secara langsung. Kedua skenario memakai algoritma \textit{Inductive Miner infrequent} (IMf) dengan \textit{noise threshold} 0,2 melalui pustaka PM4Py \cite{leemans2013discovering, berti2019pm4py}. IMf dipilih karena \textit{event log} pemeliharaan memuat perilaku jarang dan variasi label yang tinggi; pembuangan busur DFG di bawah ambang 0,2 dirancang untuk mengurangi \textit{noise} tanpa menghilangkan pola utama proses \cite{leemans2013discovering, augusto2019automated}.

\subsection{Skenario As-Is (Baseline)}

Skenario \textit{as-is} membangun model langsung dari \textit{raw labels}, yaitu kolom \texttt{activity} pada \textit{event log as-is}. Skenario ini menggambarkan kondisi proses apa adanya tanpa anotasi LLM, dan menjadi acuan perbandingan.

\subsection{Skenario Enriched AI}

Skenario \textit{enriched} membangun model dari label kategori hasil anotasi, yaitu kolom \texttt{action\_type} pada \textit{event log enriched}. Skenario ini diharapkan menghasilkan model yang lebih ringkas karena ratusan \textit{raw labels} telah diabstraksi menjadi dua belas kategori baku.

\section{Perancangan Evaluasi Model}

Kualitas model dievaluasi melalui empat dimensi \textit{conformance} yang lazim dalam \textit{process mining} \cite{buijs2014quality, vanderAalst2016_process_mining}. \textit{Fitness} mengukur sejauh mana model mampu melakukan \textit{replay} pada \textit{traces} yang terdapat di dalam \textit{event log}, dihitung melalui \textit{token-based replay}. \textit{Precision} mengukur seberapa ketat model membatasi perilaku yang tidak terjadi pada data, dihitung melalui pendekatan ETConformance. \textit{Generalization} menilai kemampuan model menangani perilaku wajar yang belum terlihat pada data. \textit{Simplicity} menilai kesederhanaan struktur model berdasarkan derajat keterhubungan \textit{Petri net}. Nilai aktual keempat metrik untuk kedua skenario dilaporkan pada Bab~IV.
